\section{Introduccción}

El presente trabajo práctico tiene como objetivo familiarizarse con los instrumentos utilizados para medir magnitudes en circuitos de corriente alterna, como el vatímetro, el amperímetro y el voltímetro. También se aplicarán los conocimientos teóricos de corriente alterna de un circuito real, estudiando la relación entre la potencia activa, reactiva y aparente. Finalmente, se analizará el método de corrección del factor de potencia mediante bancos de capacitores y el efecto producido por la incorporación de un núcleo en un inductor conectado al circuito.

El circuito utilizado en la experiencia consiste en una fuente de tensión alterna variable, una carga resistiva formada por lámparas incandescentes, representada por una resistencia $R$, en nuestro caso se utilizaron dos lámparas encendidas, una bobina modelada mediante una inductancia $L$ y su resistencia serie asociada $R_L$, y un banco de capacitores conectable en paralelo. El esquemático correspondiente se muestra en la Fig.~\ref{fig:circuito}. La tensión eficaz aplicada sobre la carga se denota mediante el fasor $V$, mientras que la corriente eficaz consumida se representa mediante el fasor $I$. En las expresiones de potencia se utilizarán los módulos de los fasores.

En corriente alterna, la potencia depende de los valores eficaces de tensión y corriente y también del desfase existente entre ambas magnitudes. La potencia aparente $S$ se define como:

\[S = |V||I| \quad (1)\]

mientras que la potencia activa $P$ viene dada por:

\[P = |V||I|\cos(\varphi) \quad (2)\]

donde $\varphi$ representa el ángulo de desfase entre la tensión y la corriente. La potencia reactiva $Q$ se expresa como:

\[Q = |V||I|\sin(\varphi) \quad (3)\]

Estas magnitudes se relacionan mediante el triángulo de potencias:

\[S^2 = P^2 + Q^2 \quad (4)\]

El factor de potencia se define como:

\[fp = \cos(\varphi) = \frac{P}{S} \quad (5)\]

La compensación del factor de potencia se realizó mediante capacitores conectados en paralelo, con el objetivo de reducir la potencia reactiva inductiva del circuito.

\begin{figure}[H]
\centering

\begin{circuitikz}

% Fuente
\draw
(0,0) to[sV,l=$V_S$] (0,4);

% Amperímetro
\draw (2,4) circle (0.5);
\node at (2,4) {A};

\draw (0,4) -- (1.5,4);
\draw (2.5,4) -- (3,4);

% Vatímetro
\draw (4,4) circle (0.5);
\node at (4,4) {W};

\draw (3,4) -- (3.5,4);
\draw (4.5,4) -- (10,4);

% Conexión vatímetro inferior
\draw (4,3.5) -- (4,0);

% Conexión vatímetro superior
\draw (4,4.5) -- (5,4.5) -- (5,4);

% Línea inferior
\draw (0,0) -- (10,0);

% Voltímetro
\draw (5.5,2) circle (0.5);
\node at (5.5,2) {V};

\draw (5.5,4) -- (5.5,2.5);
\draw (5.5,1.5) -- (5.5,0);

% Rama resistiva
\draw
(6.5,4) to[closing switch,l=$S_R$] (6.5,3.2)
to[R,l=$R$] (6.5,0);

% Rama inductiva
\draw
(8,4) to[closing switch,l=$S_L$] (8,3.2)
to[R,l=$R_L$] (8,2)
to[L,l=$L$] (8,0);

% Rama capacitiva
\draw
(10,4) to[closing switch,l=$S_C$] (10,3.2)
to[C,l=$C$] (10,0);

\end{circuitikz}

\caption{Circuito de análisis.}
\label{fig:circuito}

\end{figure}

\newpage
